% GENERATED FILE. Edit canonical lessons, then run npm run generate:books.
% canonical-source-hash: fnv1a64:ee4ce597259a15f4
% canonical-lessons: GE-C176-bedeuten, GE-C176-beeilen, GE-C176-behalten, GE-C176-beschreiben, GE-C176-betreten

\chapter{To mean, To hurry, To keep, To describe, To enter}
\label{ch:ge-to-mean-to-hurry-to-keep-to-describe-to-enter}
\hlchaptermodality{176}

\begin{chapteropening}
\textbf{By the end of this chapter:} \emph{I can say to mean, to hurry, to keep, to describe and to enter.}

Complete the last lesson of chapter 176: i can say to mean, to hurry, to keep, to describe and to enter.
\end{chapteropening}

\section[bedeuten]{bedeuten — to mean}

\label{lesson:GE-C176-bedeuten}



\begin{quote}
Before the new one: say the German for to bathe, then the German for to build.
\end{quote}

\subsection*{You'll want to know: bedeuten}
\textbf{bedeuten} — \textquotedblleft{}to mean\textquotedblright{}.

\begin{grammarlens}[title={What you've built}]
One more word for this chapter.
\end{grammarlens}

\subsection*{Your turn}
\begin{itemize}
\raggedright
  \item \emph{Say it:} \emph{bedeuten}
  \item \emph{Say it:} \emph{bedeuten}, once more
  \item \emph{Say it:} say \emph{bedeuten}
  \item \emph{Recall:} say the German for to bathe, then the German for to build, then say \emph{bedeuten} again
\end{itemize}

\begin{tcolorbox}[breakable,colback=teal!4,colframe=teal!35!black,title={Before you move on}]
What does bedeuten mean? (\textquotedblleft{}To mean\textquotedblright{}.) Say it once more.
\end{tcolorbox}
\section[beeilen]{beeilen — to hurry (sich beeilen)}

\label{lesson:GE-C176-beeilen}



\begin{quote}
Before the new one: say the German for to build, then the German for to mean.
\end{quote}

\subsection*{You'll want to know: beeilen}
\textbf{beeilen} — \textquotedblleft{}to hurry (sich beeilen)\textquotedblright{}.

\begin{grammarlens}[title={What you've built}]
One more word for this chapter.
\end{grammarlens}

\subsection*{Your turn}
\begin{itemize}
\raggedright
  \item \emph{Say it:} \emph{beeilen}
  \item \emph{Say it:} \emph{beeilen}, once more
  \item \emph{Say it:} say \emph{beeilen}
  \item \emph{Recall:} say the German for to build, then the German for to mean, then say \emph{beeilen} again
\end{itemize}

\begin{tcolorbox}[breakable,colback=teal!4,colframe=teal!35!black,title={Before you move on}]
What does beeilen mean? (\textquotedblleft{}To hurry (sich beeilen)\textquotedblright{}.) Say it once more.
\end{tcolorbox}
\section[behalten]{behalten — to keep}

\label{lesson:GE-C176-behalten}



\begin{quote}
Before the new one: say the German for to mean, then the German for to hurry.
\end{quote}

\subsection*{You'll want to know: behalten}
\textbf{behalten} — \textquotedblleft{}to keep\textquotedblright{}.

\begin{grammarlens}[title={What you've built}]
One more word for this chapter.
\end{grammarlens}

\subsection*{Your turn}
\begin{itemize}
\raggedright
  \item \emph{Say it:} \emph{behalten}
  \item \emph{Say it:} \emph{behalten}, once more
  \item \emph{Say it:} say \emph{behalten}
  \item \emph{Recall:} say the German for to mean, then the German for to hurry, then say \emph{behalten} again
\end{itemize}

\begin{tcolorbox}[breakable,colback=teal!4,colframe=teal!35!black,title={Before you move on}]
What does behalten mean? (\textquotedblleft{}To keep\textquotedblright{}.) Say it once more.
\end{tcolorbox}
\section[beschreiben]{beschreiben — to describe}

\label{lesson:GE-C176-beschreiben}



\begin{quote}
Before the new one: say the German for to hurry, then the German for to keep.
\end{quote}

\subsection*{You'll want to know: beschreiben}
\textbf{beschreiben} — \textquotedblleft{}to describe\textquotedblright{}.

\begin{grammarlens}[title={What you've built}]
One more word for this chapter.
\end{grammarlens}

\subsection*{Your turn}
\begin{itemize}
\raggedright
  \item \emph{Say it:} \emph{beschreiben}
  \item \emph{Say it:} \emph{beschreiben}, once more
  \item \emph{Say it:} say \emph{beschreiben}
  \item \emph{Recall:} say the German for to hurry, then the German for to keep, then say \emph{beschreiben} again
\end{itemize}

\begin{tcolorbox}[breakable,colback=teal!4,colframe=teal!35!black,title={Before you move on}]
What does beschreiben mean? (\textquotedblleft{}To describe\textquotedblright{}.) Say it once more.
\end{tcolorbox}
\section[betreten]{betreten — to enter, to step on}

\label{lesson:GE-C176-betreten}



\begin{quote}
Before the new one: say the German for to keep, then the German for to describe.
\end{quote}

\subsection*{You'll want to know: betreten}
\textbf{betreten} — \textquotedblleft{}to enter, to step on\textquotedblright{}.

\begin{grammarlens}[title={What you've built}]
One more word for this chapter.
\end{grammarlens}

\subsection*{Your turn}
\begin{itemize}
\raggedright
  \item \emph{Say it:} \emph{betreten}
  \item \emph{Say it:} \emph{betreten}, once more
  \item \emph{Say it:} say \emph{betreten}
  \item \emph{Recall:} say the German for to keep, then the German for to describe, then say \emph{betreten} again
\end{itemize}

\begin{tcolorbox}[breakable,colback=teal!4,colframe=teal!35!black,title={Before you move on}]
What does betreten mean? (\textquotedblleft{}To enter, to step on\textquotedblright{}.) Say it once more.
\end{tcolorbox}
